\section{context/ — 上下文管理}

\subsection{目录总体作用}

为 Agent 节点提供统一上下文接口，实现 GSSC 流水线
（Gather / Select / Structure / Compress），
将对话历史、RAG 检索结果、metadata 链、长期记忆拼装为 LLM Prompt。

\subsection{文件说明}

\begin{description}[leftmargin=4em,style=nextline]
    \item[\texttt{\_\_init\_\_.py}]
    导出 \texttt{BaseContext}、\texttt{ContextManager}。

    \item[\texttt{base.py}]
    上下文抽象基类 \texttt{BaseContext}：\texttt{save()}、\texttt{load()}、\texttt{clear()}、\texttt{build()}。

    \item[\texttt{context\_manager.py}]
    默认内存上下文实现 \texttt{ContextManager}。
    \begin{itemize}
        \item \texttt{build()} 四阶段流程：
        \begin{enumerate}
            \item Gather：从 \texttt{state.messages}、\texttt{retrieved\_context}、memory 检索
            \item Select：按 \texttt{history\_mode}（full/minimal）、\texttt{context\_profile} 过滤
            \item Structure：XML 风格块（\texttt{<context type='...'>}）
            \item Compress：按 \texttt{max\_tokens} 截断
        \end{enumerate}
        \item \texttt{format\_metadata\_chain\_for\_prompt()}：格式化 metadata 链
    \end{itemize}

    \item[\texttt{context\_policy.py}]
    兼容层，重导出 \texttt{config.context\_settings} 中的策略函数：
    \texttt{should\_persona\_file\_write()}、\texttt{is\_echo\_task()} 等。
\end{description}
